\begin{englishabstract}
When deploying electromagnetic acoustic transducers on unmanned aerial vehicles for structural non-destructive testing, accurately determining the contact state between the probe and the test surface remains a fundamental challenge. Conventional methods relying on single-threshold echo amplitude detection are unreliable under flight disturbances and attitude variations. To address this core problem, this dissertation proposes a multi-modal progressive physics-constrained framework for autonomous UAV-based electromagnetic ultrasonic detection, integrating EMAT echo sequences, UAV pose dynamic residuals, and local RGBD visual ROI features into a unified temporal perception model.

At the theoretical level, this work first establishes the electromagnetic-elastic coupling mechanism, treating the EMAT signal as an observable dynamic process formed by four cascaded physical stages: excitation, propagation, reflection, and reception. The Lorentz body force is derived as the source term in the elastodynamic equation. Building on this foundation, contact boundary and lift-off effects are characterized through echo statistical modeling, defining contact state as a multi-factor latent variable jointly determined by lift-off distance, normal constraint, end-effector mechanical stability, and echo observability, expressed as a continuous probability distribution.

At the platform and data level, a multi-modal UAV experimental platform is designed and integrated. The system uses a multirotor UAV equipped with a Livox MID-360 LiDAR, Intel RealSense D435 RGBD camera, PX4 flight controller, and MAVROS communication link, with FAST-LIO 2.0 providing real-time 6-DOF state estimation. A ROS driver for the CH346C-based EMAT probe is developed, achieving stable waveform acquisition at approximately 40 Hz with Qt5 visualization. A synchronized multi-modal data acquisition system records pose, RGBD video streams, and EMAT echo signals, providing structured data for model training.

At the algorithmic level, to address the asynchronous sampling problem across EMAT (~40 Hz), flight control (~100 Hz), and RGBD vision (~30 Hz) modalities, a unified time-base weighted interpolation alignment with independent encoder mapping is proposed. Building on this, a physics-constrained attention mechanism is designed, embedding temporal proximity and spatial residual constraints as log-bias terms into the scaled dot-product attention of a Transformer, ensuring cross-modal interactions simultaneously obey temporal continuity, spatial proximity, and kinematic plausibility. Furthermore, modality reliability gating weights and a dual-threshold duration-constrained state machine are introduced, forming a complete evidence chain from sensor uncertainty assessment to control-loop hysteresis decision.

Experimental results demonstrate that the proposed multi-modal physics-constrained fusion framework effectively overcomes the limitations of single-modality threshold methods, achieving smooth state probability transitions in critical contact segments, and providing an integrated solution spanning physical modeling, signal acquisition, feature fusion, and closed-loop control for autonomous UAV-based EMAT detection.

\englishkeyword{UAV Non-Destructive Testing, Electromagnetic Acoustic Transducer, Multi-modal Fusion, Physics-Constrained Attention, Contact State Discrimination, Transformer}
\end{englishabstract}